%! AP Statistics — Unit 5, Lesson 5.7 — Lesson Plan
\documentclass[10pt]{article}
\usepackage{apstats-article}
\usepackage{apstats-boxes}

\begin{document}

\pageheader{Unit 5, Lesson 5.7}{Lesson Plan}

\begin{objectivebox}
\textbf{Topic:} 5.7 — Sampling Distributions for Sample Means\\
\textbf{Enduring Understanding:} UNC-3\\
\textbf{Learning Objectives:} UNC-3.Q, UNC-3.R, UNC-3.S\\
\textbf{Essential Knowledge:} UNC-3.Q.1; UNC-3.R.1--R.2; UNC-3.S.1--S.3\\
\textbf{Mathematical Practice:} 3.C — Verify conditions; 4.B — Connect representations
\end{objectivebox}

\begin{learningtargetbox}
By the end of this lesson, students will be able to:
\begin{itemize}
  \item State $\mu_{\bar{x}} = \mu$ and $\sigma_{\bar{x}} = \sigma/\sqrt{n}$.
  \item Check conditions (random, 10\%, Normal/Large Sample) for using the normal model.
  \item Compute probabilities for $\bar{x}$ using the normal approximation.
\end{itemize}
\end{learningtargetbox}

\begin{vocabbox}
sampling distribution of $\bar{x}$, $\mu_{\bar{x}}$, $\sigma_{\bar{x}}$, Normal/Large Sample condition ($n \ge 30$ or pop.\ normal), standard error of $\bar{x}$
\end{vocabbox}

\begin{hookbox}
\textbf{Hook — Quality Control (5 min):} A cereal company fills boxes to $\mu = 500$ g with $\sigma = 15$ g. A random sample of 36 boxes is weighed. If the sample mean falls below 495 g, the machine is recalibrated. How often does this happen by chance alone?
\end{hookbox}

\begin{notesbox}
\textbf{Direct Instruction (15 min):} State $\mu_{\bar{x}} = \mu$ (unbiased) and $\sigma_{\bar{x}} = \sigma/\sqrt{n}$. Tie Normal/Large Sample condition to CLT from Lesson 5.3. Work full probability example. Note difference: here $\sigma$ is the population SD, not sample SD.
\end{notesbox}

\begin{reflectionbox}
\textbf{Exit Ticket + Homework:} Check conditions; compute probabilities for $\bar{x}$.\\
\textbf{AP Connection:} UNC-3.Q--S are tested on most AP exams; parallel to proportion formulas.
\end{reflectionbox}

\end{document}
